\section{Approximation on polydiscs}
\label{sec:gluing}

We now pass from the planar operation $L$ to approximation in every
source dimension. The first step is to enlarge a polydisc. A change
of source coordinates will put the part beyond the original map domain
into the form of one long planar variable with bounded passive
variables. Near the joining annulus, the passive variables describe
sets of arbitrarily small diameter, so that $L$ applies to their
images in a single witness neighborhood. We then join this extension
to the original map on a buffered open overlap.

Write
\[
 \overline\Delta_R^{\,m}=\{z\in\C^m:\max_j|z_j|\leq R\},
 \qquad \|z\|_\infty=\max_j|z_j|.
\]
All holomorphic maps on compact sets below are defined on open
neighborhoods of those sets.  Extra polydiscs always have positive
radii.  A zero-dimensional polydisc is allowed.

\begin{theorem}\label{ana:cap}
Let \(Y\) be a complex surface with a complete compatible distance.
If the local operation \(L\) holds at every point of \(Y\), then \(Y\)
has the convex approximation property in every source dimension.
\end{theorem}

This section proves the polydisc case. Section~\ref{sec:globalization}
will remove the remaining restriction on the shape of the convex set.

\subsection{Gluing on a buffered open overlap}

The source change need not preserve product overlaps, so the planar
splitting of Lemma~\ref{loc:splitting} must be replaced by a splitting
on a convex neighborhood in $\C^n$. Let $Q\subset\C^n$ be compact
and convex, and put $Q_r=\{z:\operatorname{dist}(z,Q)<r\}$ for $r>0$.
An auxiliary variable $v$ ranges over a ball $B_b\subset\C^N$;
the splitting operators below act only on $z$ and preserve
holomorphic dependence on $v$.

Suppose two open sets
$A,B$ cover a neighborhood of $\overline{Q_r}$ and have a fixed buffered
overlap: there is a smooth cutoff equal to zero near $A\setminus B$ and
one near $B\setminus A$, whose transition lies inside $A\cap B$ with a
positive collar. The collar is understood relative to the total domain;
after shrinking $Q_r$ its width has a positive lower bound depending
only on the fixed geometry. For a bounded holomorphic jump $u$,
the functions $\chi u$ and $(\chi-1)u$, extended by zero where their
respective multipliers vanish, have the same $\bar\partial$ derivative.
Fix $0<h<\min(r/2,1)$. The constants below depend only on the fixed
overlap geometry, a bounded range of $r$, and the dimensions, not on
the jump $u$ or its holomorphic auxiliary parameters.
Solve this common equation on an intermediate convex thickening,
using the weighted minimal $L^2$ solution with weight $|z|^2$
\cite[Theorem~2.2.1$'$]{HormanderLTwo}. The weight is bounded above
and below on the fixed bounded domains, so the solution has $L^2$
norm at most a fixed constant times that of the input form.
Cauchy estimates on the buffered transition region control any fixed
number of derivatives of that form by powers of $h^{-1}\|u\|_\infty$.
Interior elliptic estimates, applied to
$\Delta w=4\sum_j\partial\eta_j/\partial z_j$ when
$\bar\partial w=\sum_j\eta_j\,d\bar z_j$, give the corresponding
supremum bound for $w$ on the smaller thickening. Subtracting $w$
from the two cutoff functions gives bounded linear splittings on
$A\cap Q_{r-h}$ and $B\cap Q_{r-h}$, with bound
$Ch^{-p}\|u\|_\infty$ for some fixed integer $p$.
The minimal solution operator is linear and fixed on the base domain;
it therefore preserves holomorphic auxiliary parameters.

\begin{lemma}[Small nonlinear gluing on an open overlap]\label{glue:open}
In the preceding buffered open-overlap geometry, let two holomorphic
sprays into a complex manifold be related on their overlap by a fiber
coordinate change
\[
 \gamma(z,v)=v+c(z,v),\qquad v\in B_b.
\]
Suppose $\|c\|_\infty$ can be made arbitrarily small with $b>0$ fixed.
After any prescribed positive loss of base and fiber margins, the two
sprays admit fiber reparametrizations close to the identity which make
them agree. Evaluating at $v=0$ gives a glued holomorphic map. It is
arbitrarily close to the original map on each compactum where the
corresponding input spray is fixed, or has a uniform modulus of
continuity in the fiber variable. No such uniform target estimate is
asserted on a side whose spray varies without this control.
\end{lemma}

\begin{proof}
Use the additive splitting to find $a,b$ with $b-a=c$. Shrink the
base and fiber domains by $h$. The splitting bounds and Cauchy
estimates in $v$ imply, with a fixed integer $p$ enlarged as necessary,
\[
 \|a\|+\|b\|\le Ch^{-p}e,\qquad e=\|c\|,
\]
including their needed fiber derivatives on intermediate smaller balls.
For $Ch^{-p}e$ sufficiently small relative to these margins,
$\operatorname{id}+b$ has a fiber inverse on a smaller ball, uniformly
in the base. The new transition is
\[
 (\operatorname{id}+b)^{-1}\circ\gamma\circ(\operatorname{id}+a)
       =\operatorname{id}+d.
\]
The equation determining $d$ is
\[
 d= c(v+a)-c(v)-\{b(v+d)-b(v)\}.
\]
The first difference is quadratic in $e$, up to powers of $h^{-1}$.
The operation in braces is contractive on a small supremum-norm ball,
by the first fiber derivative bound. It follows that
\begin{equation}\label{eq:glue-quadratic}
             \|d\|\le Ch^{-p}e^2.
\end{equation}
These estimates involve a fixed finite number of derivatives. Thus one
common exponent $p$ works throughout the iteration; no derivatives of
increasing order are requested.

Choose losses $h_i=h_0 2^{-i}$ whose sum fits within the prescribed
margins, leaving intermediate buffers at each step. With
$A=Ch_0^{-p}$, Equation~\eqref{eq:glue-quadratic} reads
$e_{i+1}\le A2^{pi}e_i^2$. Put $E_i=A2^{p(i+1)}e_i$. Then
$E_{i+1}\le E_i^2$. Choosing $E_0<1/2$ gives doubly exponential decay.
In particular all sums $\sum_i h_i^{-q}e_i$ converge for every fixed
$q$, and are as small as desired with $e_0$.

The successive fiber reparametrizations are composed on the nested
domains. Their displacements fit within the reserved buffers and their
fiber derivatives have summable deviations from the identity. The
compositions therefore converge uniformly. Their limiting transitions
are the identity.
Composing the two original sprays with these limiting side maps gives
equal maps on the overlap.
The target approximation assertion follows from the stated continuity
control and the smallness of the limiting fiber corrections.
\end{proof}

\subsection{A change of coordinates with a thin parameter range}

Enlarging one coordinate of a polydisc directly leaves an entire
transverse polydisc over each boundary point; its image need not fit
one witness neighborhood for $L$. We instead deform the source by a
polynomial automorphism. It will be nearly the identity on the old
polydisc, while the rest of the enlarged source has coordinates
$(s,\xi)$ with $s$ planar and $\xi$ bounded. The following estimate
constructs these coordinates on inverse images of the unit polydisc.
The factors $d^{-1}$ in their formula will then make the $x$-image
thin as $\xi$ varies with $s$ fixed on the joining annulus.

\begin{lemma}\label{ana:concentration}
Fix \(m\geq2\), numbers
\[
 1<r_m<\cdots<r_2<r_1,
 \qquad 0<a<1,\qquad r_j<a r_{j-1}\quad(2\leq j\leq m).
\]
For an integer \(d\geq2\), set
\[
 P_j(w)=(w/r_j)^d,\qquad
 H_d(x)=\bigl(a x_m,P_2(x_2)-a x_1,\ldots,
                        P_m(x_m)-a x_{m-1}\bigr).
\]
There are constants \(C\geq1\), \(q\in(0,1)\), and \(d_0\), depending
only on the displayed fixed data, with the following properties.
For \(d\geq d_0\), \(N\geq1\), and
\[
 E_{d,N}=H_d^{-N}(\overline\Delta_1^{\,m}),
\]
one has
\begin{equation}\label{ana:eq:power-bound}
 \max_{j\geq2}|P_j(x_j)|\leq C\max(1,|x_1|)
 \qquad(x\in E_{d,N}).
\end{equation}
On \(E_{d,N}\cap\{|x_1|>r_1\}\), all \(x_j\) are nonzero and
\begin{equation}\label{ana:eq:ratio-bound}
 \left|\frac{H_d(x)_j}{a x_{j-1}}\right|\leq q
 \qquad(2\leq j\leq m).
\end{equation}
Consequently the holomorphic functions
\[
 \xi_j(x)=\log\left(\frac{P_j(x_j)}{a x_{j-1}}\right),
 \qquad 2\leq j\leq m,
\]
using the logarithm on \(\{|\zeta-1|<1\}\), satisfy
\[
 |\xi_j(x)|\leq B:=-\log(1-q).
\]
Here the assertions about holomorphicity are on any open inverse
image of the unit polydisc; all the estimates also hold on its
closure.  With
\[
 D=d^{m-1},\qquad
 \rho_1=1,\qquad \rho_j=r_j(a\rho_{j-1})^{1/d},
\]
there is a single-valued holomorphic function \(s\) on this end such
that
\begin{equation}\label{ana:eq:parameterization}
 x_j=\rho_j s^{d^{m-j}}
       \exp\left(\sum_{k=2}^j\frac{\xi_k}{d^{j-k+1}}\right),
 \qquad 1\leq j\leq m.
\end{equation}
For \(j=1\), the sum is empty and this says \(x_1=s^D\).
\end{lemma}

\begin{proof}
Put \(r_*=r_1\).  Choose \(S>\max(2,r_*)\), and then \(d_0\geq2\)
large enough that
\[
 (M/r_*)^d\geq4M\qquad(M\geq S,\ d\geq d_0).
\]
For example, it suffices to impose this at \(M=S,d=d_0\), since
\((M/r_*)^d/M\) increases in both variables on this range.
The cone
\[
 {\cal E}=\{u:\max_{j\geq2}|u_j|>\max(S,|u_1|)\}
\]
is forward invariant under \(H_d\).  Indeed, if
\(M=\max_{j\geq2}|u_j|\) and \(|u_k|=M\), then
\[
 |H_d(u)_k|\geq(M/r_*)^d-aM\geq3M,
 \qquad |H_d(u)_1|\leq aM.
\]
The transverse maximum therefore grows by at least a factor \(3\).
In particular, no forward orbit that ends in
\(\overline\Delta_1^{\,m}\) can enter \({\cal E}\).

Choose \(C\) larger than
\[
 1,\quad 4a,\quad 16a^2r_*^2,\quad 4S.
\]
If \(A=\max_{j\geq2}|P_j(x_j)|>C\max(1,|x_1|)\), then
\[
 \max_{j\geq2}|x_j|\leq r_* A^{1/d}\leq r_*\sqrt A.
\]
For an index \(k\) attaining \(A\), both possible terms
\(a|x_{k-1}|\), namely \(a|x_1|\) when \(k=2\) and a transverse
coordinate otherwise, are at most \(A/4\).  Hence
\[
 |H_d(x)_k|\geq3A/4,\qquad
 |H_d(x)_1|\leq a r_*\sqrt A\leq A/4.
\]
Thus \(H_d(x)\in{\cal E}\), a contradiction.  This proves
\eqref{ana:eq:power-bound}; the same argument applies at every
nonfinal point of the orbit.  At its final point the bound is
immediate.

Write \(y=H_d(x)\).  Applying the power bound at \(y\), or using
\(\|y\|_\infty\leq1\) if \(y\) is the final point, gives
\begin{equation}\label{ana:eq:next-coordinate}
 |y_j|\leq r_j\bigl(C\max(1,a|x_m|)\bigr)^{1/d}
 \qquad(j\geq2).
\end{equation}
Let \(L_j=\log^+|x_j|\), \(M=L_m\), and
\(A_1=\log(r_*\sqrt C+a)\).  From
\(P_j(x_j)=y_j+a x_{j-1}\) and
\eqref{ana:eq:next-coordinate},
\[
 |P_j(x_j)|
 \leq (r_*\sqrt C+a)\exp(L_{j-1}+M/d).
\]
Whether or not \(|x_j|\) exceeds \(1\), this implies
\begin{equation}\label{ana:eq:log-recursion}
 L_j\leq \log r_j+A_1/d+L_{j-1}/d+M/d^2.
\end{equation}
Put \(A_2=\log r_*+A_1/2\).  Iterating
\eqref{ana:eq:log-recursion} to \(j=m\) yields
\[
 M\leq 2A_2+d^{-(m-1)}L_1+
       M\sum_{i=0}^{m-2}d^{-2-i}
 \leq 2A_2+d^{-(m-1)}L_1+\tfrac12M.
\]
For \(|x_1|\geq r_1\), define
\[
 t=(|x_1|/r_1)^{1/D}\geq1,\qquad
 A_3=4A_2+2\log r_1.
\]
It follows that
\[
 M\leq A_3+2\log t,\qquad
 |y_j|\leq r_j\exp((\log C+A_3)/d)t^{2/d}.
\]

Choose a sufficiently small \(\delta>0\) so that
\[
 q:=\max_{2\leq j\leq m}
       \frac{r_j+\delta}{a(r_{j-1}-\delta)}<1,
 \qquad r_j-\delta>0.
\]
Increase \(d_0\), if necessary, so that for all \(d\geq d_0\)
\[
 r_j\exp((\log C+A_3)/d)\leq r_j+\delta,\qquad
 r_j\bigl((1-q)a(r_{j-1}-\delta)\bigr)^{1/d}
       \geq r_j-\delta .
\]
We prove successively that
\[
 |x_j|\geq(r_j-\delta)t^{d^{m-j}}\quad(1\leq j\leq m).
\]
The assertion for \(j=1\) follows from the definition of \(t\).
If it holds for \(j-1\), then
\[
 \frac{|y_j|}{a|x_{j-1}|}
 \leq
 \frac{r_j+\delta}{a(r_{j-1}-\delta)}
 t^{\,2/d-d^{m-j+1}}\leq q.
\]
Therefore
\[
 |P_j(x_j)|\geq(1-q)a|x_{j-1}|,
\]
and taking the \(d\)-th root proves the next assertion.
This also proves \eqref{ana:eq:ratio-bound}.
Since \(P_j/(a x_{j-1})=1+y_j/(a x_{j-1})\), the indicated
logarithm is well defined; its power series gives the bound
\(|\xi_j|\leq\sum_{\nu\geq1}q^\nu/\nu=B\).

There is no root choice in the definition of \(s\).  Define it by
\[
 s=\frac{x_m}{\rho_m}
       \exp\left(-\sum_{k=2}^m\frac{\xi_k}{d^{m-k+1}}\right).
\]
The equations
\[
 x_j^d=r_j^d a x_{j-1}\exp(\xi_j)
\]
then prove \eqref{ana:eq:parameterization} backwards from \(j=m\)
to \(j=1\), using the definition of the positive real numbers
\(\rho_j\).
\end{proof}

\begin{proposition}\label{ana:polydisc}
Under the hypotheses of Theorem~\ref{ana:cap}, a map holomorphic near
a closed polydisc in \(\C^m\) can be approximated uniformly there by
entire maps \(\C^m\to Y\).  More precisely, a map near
\(\overline\Delta_1^{\,m}\) can be approximated there by maps
holomorphic near \(\overline\Delta_R^{\,m}\), for every finite
\(R>1\).
\end{proposition}

\begin{proof}
First prove the finite-radius assertion.  When \(m=1\), the
one-dimensional boundary circle is covered by finitely many of
the neighborhoods in \(L\), and Proposition~\ref{loc:planar}
applies with no passive variables.

Suppose \(m\geq2\).  Let \(f\) be the initial map and let
\(\varepsilon>0\).  Choose \(b>1\) with \(f\) holomorphic near
\(\overline\Delta_b^{\,m}\).  Choose the radii \(r_j\), the number
\(a\), and \(\ell>0\) so that Lemma~\ref{ana:concentration} applies
and
\[
 1<r_m<\cdots<r_1<r_1+3\ell<b.
\]
The derivative of \(H_d\) at zero is the \(d\)-independent map
\[
 L(x)=(a x_m,-a x_1,\ldots,-a x_{m-1}),
 \qquad \|Lx\|_\infty=a\|x\|_\infty.
\]
The map \(H_d\) is an automorphism: from \(y=H_d(x)\), first recover
\(x_m=y_1/a\), and then successively recover
\[
 x_{j-1}=(P_j(x_j)-y_j)/a,\qquad j=m,m-1,\ldots,2.
\]
Fix \(R_+>R\), and fix \(N\geq1\) so that \(a^N R_+<1\).
Set
\[
 h_d=H_d^{-N}L^N,\qquad \Omega=\Delta_{R_+}^{\,m}.
\]
Thus \(h_d(\overline\Omega)\subset E_{d,N}\).

On some fixed neighborhood of \(\overline\Delta_1^{\,m}\),
\begin{equation}\label{ana:eq:near-identity}
 h_d\longrightarrow\operatorname{id}\quad(d\longrightarrow\infty).
\end{equation}
Here is a quantitative justification.  Choose
\(1<b_0<b_1<r_m\), put
\(\eta=(b_1-b_0)/3\), and put \(\theta=b_1/r_m<1\).
If \(\|L^{-1}y\|_\infty\leq b_0+\eta\) and
\(\theta^d/a<\eta\), the inverse recursion above gives
\[
 \|H_d^{-1}(y)-L^{-1}(y)\|_\infty\leq\theta^d/a.
\]
Indeed, its first recovered coordinate is exact, and each next
one differs from the corresponding coordinate of \(L^{-1}y\)
by \(P_j(x_j)/a\).  Inductively all recovered coordinates have
modulus at most \(b_0+2\eta<b_1\), so that
\(|P_j(x_j)|\leq\theta^d\).
Let \(C_N=\sum_{i=1}^N a^{-i}\), and take \(d\) large enough
that \(C_N\theta^d/a<\eta\).  Successive comparison of
\(H_d^{-r}L^N(z)\) and \(L^{N-r}(z)\), for
\(z\in\overline\Delta_{b_0}^{\,m}\), gives the error recursion
\[
 e_0=0,\qquad e_{r+1}\leq a^{-1}e_r+\theta^d/a,\qquad
 e_r\leq\theta^d\sum_{i=1}^r a^{-i}.
\]
The imposed bound guarantees the hypothesis of the inverse
estimate at every step.  At \(r=N\) it proves
\(\|h_d-\operatorname{id}\|_{\overline\Delta_{b_0}^{\,m}}
\leq C_N\theta^d\), and hence \eqref{ana:eq:near-identity}.

Choose a number \(b'<b\) larger than \(r_1+2\ell\).
Lemma~\ref{loc:spray} gives a chart spray
\[
 f^\sharp(x,v),\qquad f^\sharp(x,0)=f(x),
\]
on a neighborhood of
\(\overline\Delta_{b'}^{\,m}\times\overline\Delta_\beta^{\,2}\)
for some \(\beta>0\).  Its vertical differential is invertible and
its vertical maps are injective after reducing \(\beta\).
The value \(b'\) and this spray are fixed before \(d\) is increased.

Write \(\Phi_d(s,\xi)\) for the right side of
\eqref{ana:eq:parameterization}, now allowing \(\xi\) to vary
independently.  Use, for example, the three parameter polydiscs
\[
 \Xi_i=\{\xi\in\C^{m-1}:|\xi_j|\leq B+i\},\qquad i=1,2,3.
\]
Choose fixed numbers
\[
 r_1<c_0<c_1<c_2<c_3<r_1+\ell.
\]
On
\[
 c_0\leq |s|^D\leq c_3,\qquad \xi\in\Xi_3,
\]
one has, uniformly as \(d\to\infty\),
\begin{equation}\label{ana:eq:thin-range}
 |\Phi_{d,j}(s,\xi)|\leq r_j\exp(A_4/d)\quad(j\geq2),
 \qquad
 \sup_{\xi,\widetilde\xi\in\Xi_3}
 \|\Phi_d(s,\xi)-\Phi_d(s,\widetilde\xi)\|_\infty
       \leq A_5/d,
\end{equation}
where the following fixed constants suffice:
\[
 \begin{aligned}
 R_0&=2\log r_1+|\log a|,\\
 A_4&=|\log a|+R_0+\log c_3+2(B+3),\\
 A_5&=4(B+3)r_1e^{A_4/2}.
 \end{aligned}
\]
Indeed the recurrence for \(\rho_j\) gives
\[
 |\log\rho_j|\leq R_0,\qquad
 |\log\rho_j-\log r_j|
       \leq (|\log a|+R_0)/d .
\]
The power of \(s\) contributes at most
\(\log c_3/d^{j-1}\leq\log c_3/d\).  The remaining exponential
satisfies
\[
 \sum_{k=2}^j|\xi_k|/d^{j-k+1}\leq2(B+3)/d .
\]
Moreover
\[
 \partial_{\xi_k}\Phi_{d,j}
       =d^{-(j-k+1)}\Phi_{d,j}\quad(k\leq j),\qquad
 \partial_{\xi_k}\Phi_{d,j}=0\quad(k>j),
\]
Integrating along line segments in the convex polydisc \(\Xi_3\)
and using \(\sum_{\nu\geq1}d^{-\nu}\leq2/d\) proves the diameter
bound with the displayed \(A_5\).
In particular \(\Phi_d\) takes the indicated product into
\(\Delta_{b'}^{\,m}\) for large \(d\).

The compact set \(f(\overline\Delta_{b'}^{\,m})\) has a finite cover
by neighborhoods on which \(L\) is available.  A Lebesgue number
for a slightly smaller such cover, uniform continuity of \(f\),
and \eqref{ana:eq:thin-range} show the following.  After increasing
\(d\) and decreasing the fixed spray radius, all the values
\[
 f^\sharp(\Phi_d(s,\xi),v),\qquad \xi\in\Xi_3,\quad
 v\in\overline\Delta_{\beta}^{\,2},
\]
belong to one of these neighborhoods for each fixed \(s\) on the
outer circle \(|s|^D=c_3\).  The chosen neighborhood may depend
on \(s\).  The same increase of \(d\) ensures that
\[
 \sup_{\overline\Delta_1^{\,m}}d_Y(f(h_d(z)),f(z))
       <\varepsilon/2,\qquad
 |h_d(z)_1|<c_1\quad(z\in\overline\Delta_1^{\,m}).
\]
Require also \(r_j(Cc_2)^{1/d}<b'\) for \(j\geq2\); this
is possible since \(r_j<b'\).
Fix this value of \(d\), henceforth writing \(h=h_d\).

Apply Proposition~\ref{loc:planar} to the annulus
\[
 A_s=\{c_0\leq |s|^D\leq c_3\}
\]
and to the map \(f^\sharp(\Phi_d(s,\xi),v)\), with \(\xi,v\)
as passive variables.  It gives arbitrarily close approximation
on \(A_s\times\Xi_2\times\overline\Delta_{\beta/2}^{\,2}\)
by a holomorphic map \(G(s,\xi,v)\) defined out to any prescribed
finite outer radius in \(s\), with a smaller fixed positive
parameter margin.  Choose that radius larger than \(c_3^{1/D}\)
and larger than
\[
 \sup\{|s(h(z))|:z\in\overline\Omega,\ |h(z)_1|\geq c_1\}.
\]
If the set is empty the latter requirement is vacuous.
Otherwise this supremum is finite by compactness and the explicit formula
for \(s\); \(|s(h(z))|^D=|h(z)_1|\).
The inner radius of the domain of \(G\) can be kept strictly below
\(c_1^{1/D}\).

On the two open subsets of \(\Omega\)
\[
 A=\{z:|h(z)_1|<c_2\},\qquad
 B'=\{z:|h(z)_1|>c_1\},
\]
consider the maps with the common auxiliary variable \(v\)
\[
 F_A(z,v)=f^\sharp(h(z),v),\qquad
 F_{B'}(z,v)=G(s(h(z)),\xi(h(z)),v).
\]
The first is defined on all of \(A\): by
\eqref{ana:eq:power-bound}, its transverse coordinates satisfy
\[
 |h(z)_j|\leq r_j(Cc_2)^{1/d}<b'
\]
after the choices already made.  The second is defined on all
of \(B'\) by Lemma~\ref{ana:concentration} and the choices of
the parameter and outer-radius margins.
On their overlap, \(\Phi_d(s(h(z)),\xi(h(z)))=h(z)\);
therefore \(F_{B'}\) approximates \(F_A\) arbitrarily closely,
uniformly for \(v\) in a fixed smaller polydisc.

Choose a smooth cutoff of \(|h(z)_1|^2\) whose transition is
compactly contained in \(c_1<|h(z)_1|<c_2\).
On the bounded domain \(\Omega\), all its derivatives needed in
Lemma~\ref{glue:open} are bounded.  These constants, \(d,N\),
and all the domain margins have been fixed before choosing the
accuracy in Proposition~\ref{loc:planar}.  Lemma~\ref{glue:open}
therefore glues the two maps, by arbitrarily small vertical
corrections, on a neighborhood of \(\overline\Delta_R^{\,m}\).
On \(\overline\Delta_1^{\,m}\subset A\) the resulting central
map differs from \(f\circ h\) by less than \(\varepsilon/2\).
Together with \eqref{ana:eq:near-identity}, this proves the
finite-radius assertion.

To obtain an entire map, use this assertion, after affine
rescaling, successively on exhausting polydiscs.  Choose positive
errors with sum less than the prescribed final error, and ensure
that the error at each stage is controlled on the previous
polydisc.  Completeness makes the maps uniformly Cauchy on every
compact set.  Their limit is holomorphic: near each point, choose
a coordinate neighborhood of the limit value; uniform convergence
puts all sufficiently late maps in that chart on a smaller source
neighborhood, where the usual theorem for uniformly convergent
holomorphic functions applies.  The same argument gives the
prescribed approximation on the initial polydisc.
\end{proof}
